% !TEX root = ../../tree_notes.tex
% =====================================================================
%  Chapter 3 opener -- "Chapter at a Glance", roadmap, topic table,
%  notation. NO numbered boxes here (they would be numbered 3.0.x).
% =====================================================================

\begin{chaptertheorybox}{Chapter at a Glance}
A tree stores its nodes in \emph{space}, not in a line. To use a tree you first have to
\term{traverse} it: walk over every node, \emph{visiting} each one exactly once in some
order. "Visit" just means \emph{do your work on the node} -- print its value, add it to a
sum, check a condition. The only real choice is \textbf{when} you visit a node relative to
its subtrees, and that one choice gives the four orders this chapter is about:

\begin{itemize}[nosep]
  \item \term{preorder}: visit the node \emph{before} its two subtrees;
  \item \term{inorder}: visit it \emph{between} the left subtree and the right subtree;
  \item \term{postorder}: visit it \emph{after} both subtrees are done;
  \item \term{level order}: forget subtrees -- go level by level, top to bottom, using a
        queue instead of recursion.
\end{itemize}

\textbf{You will learn to}
\begin{itemize}[nosep]
  \item produce all four orders by hand on any small tree, and know \emph{why} each order
        is the natural choice for a different job (copy $=$ preorder, delete $=$ postorder,
        level-by-level $=$ BFS);
  \item write each order recursively in three lines, then again with an explicit stack,
        and again with Morris threading in $O(1)$ extra memory;
  \item split a level-order walk into one vector per level -- the template behind
        LeetCode 102/103/107/199/637;
  \item rebuild a whole tree from two of its traversal lists, and flatten a tree to a
        string and back (serialise/deserialise).
\end{itemize}

\textbf{Prerequisites.} \chref{c02} -- the \code{TreeNode} struct (\defref{c02-treenode})
and the structural-recursion idea (\secref{c02-model}). We also reuse the null-link count
of \thmref{c02-null-links} when we prove Morris traversal is linear.

\textbf{In contests and interviews.} Traversal questions are the most-asked tree family on
LeetCode. The recursive answers are short; the follow-ups -- "do it iteratively", "do it in
$O(1)$ space", "rebuild the tree from these two lists" -- are what separate a passing from
a strong interview, and \secref{c03-iterative} and \secref{c03-rebuild} train exactly those.
\end{chaptertheorybox}

% ---------------- The running example, introduced early ----------------
\subsection*{The Tree We Will Use All Chapter}

Every trace in this chapter runs on the \emph{same} six-node tree, so you can always come
back to one picture instead of learning a new shape per section:

\begin{figure}[H]
\centering
\begin{tikzpicture}[diagram]
  \node[troot] (r)  at (0,0)     {$1$};
  \node[tnode] (l)  at (-2,-1.4) {$2$};
  \node[tnode] (c)  at (2,-1.4)  {$3$};
  \node[tleaf] (ll) at (-2.8,-2.8) {$4$};
  \node[tleaf] (lr) at (-1.2,-2.8) {$5$};
  \node[nullnode] (nl) at (1.2,-2.8) {};
  \node[tleaf] (rr) at (2.8,-2.8) {$6$};
  \draw[tedge] (r) -- (l); \draw[tedge] (r) -- (c);
  \draw[tedge] (l) -- (ll); \draw[tedge] (l) -- (lr);
  \draw[guide] (c) -- (nl);
  \draw[tedge] (c) -- (rr);
\end{tikzpicture}
\caption{The running example $T$. Node $3$ has \emph{no left child} -- deliberately kept,
because one-child nodes are exactly where traversal mistakes show up. Its four orders:
preorder $1,2,4,5,3,6$; inorder $4,2,5,1,3,6$; postorder $4,5,2,6,3,1$;
level order $1,2,3,4,5,6$. We derive every one of these, step by step, in the chapter.}\label{fig:c03-runner}
\end{figure}

% ---------------- Roadmap diagram ----------------
\begin{figure}[H]
\centering
\begin{tikzpicture}[diagram, node distance=4mm,
    top/.style={boxnode, fill=primaryblue!10, draw=primaryblue, text width=28mm,
                minimum height=10mm, font=\small},
    res/.style={annot, text width=28mm, anchor=north}]
  \node[top] (t1) {\textbf{3.1} The three DFS orders};
  \node[top, right=of t1] (t2) {\textbf{3.2} Iterative \& Morris};
  \node[top, right=of t2] (t3) {\textbf{3.3} Level order (BFS)};
  \node[top, right=of t3] (t4) {\textbf{3.4} Rebuild \& serialise};
  \draw[flow] (t1) -- (t2); \draw[flow] (t2) -- (t3); \draw[flow] (t3) -- (t4);
  \node[res] at ([yshift=-2mm]t1.south) {pre/in/post by recursion; which order for
    which job};
  \node[res] at ([yshift=-2mm]t2.south) {same orders with your own stack; $O(1)$ space
    Morris};
  \node[res] at ([yshift=-2mm]t3.south) {queue; split into levels; zigzag, bottom-up};
  \node[res] at ([yshift=-2mm]t4.south) {pre$+$in rebuild; encode to string; Euler tour
    teaser};
\end{tikzpicture}
\caption*{The four topics of this chapter. 3.1 defines the orders, 3.2 re-implements them
without recursion, 3.3 adds the fourth (queue-based) order, and 3.4 turns the whole idea
around: from orders back to trees.}
\end{figure}

% ---------------- Topic table ----------------
\begin{center}
\renewcommand{\arraystretch}{1.25}
\begin{tabular}{@{}p{0.34\linewidth} p{0.38\linewidth} p{0.20\linewidth}@{}}
\toprule
\textbf{Topic} & \textbf{Key results} & \textbf{Solved problems} \\
\midrule
\secref{c03-dfs}\ \ The three depth-first orders
  & recursive definitions; root position in each order
  & \probref{c03-preorder}, \probref{c03-inorder}, \probref{c03-postorder} \\
\addlinespace
\secref{c03-iterative}\ \ Iterative traversals \& Morris
  & explicit-stack simulations; $O(1)$-space inorder
  & \probref{c03-nary-pre}, \probref{c03-nary-post} \\
\addlinespace
\secref{c03-bfs}\ \ Breadth-first / level order
  & queue with level splits; BFS assigns depths
  & \probref{c03-level-bottom}, \probref{c03-zigzag} \\
\addlinespace
\secref{c03-rebuild}\ \ Rebuilding and serialising
  & pre $+$ in $\Rightarrow$ unique tree; string round-trip
  & \probref{c03-build-prein}, \probref{c03-build-inpost}, \probref{c03-serialize} \\
\bottomrule
\end{tabular}
\end{center}

% ---------------- Notation ----------------
\subsection*{Notation Used in This Chapter}

\begin{center}
\renewcommand{\arraystretch}{1.2}
\begin{tabular}{@{}ll@{\qquad}ll@{}}
\toprule
\textbf{Symbol} & \textbf{Meaning} & \textbf{Symbol} & \textbf{Meaning} \\
\midrule
$\text{pre}(T)$ & preorder list of $T$   & $\text{in}(T)$  & inorder list of $T$ \\
$\text{post}(T)$ & postorder list of $T$ & $\text{lev}(T)$ & level-order list of $T$ \\
$S$             & the explicit stack (top at left) & $Q$ & the FIFO queue (front at left) \\
$\tin(v)$       & entry time of $v$ in a DFS walk  & $\tout(v)$ & exit time of $v$ \\
$v\cdot X$      & value of $v$ followed by list $X$ & \code{\#} & null marker in strings \\
\bottomrule
\end{tabular}
\end{center}

Throughout, \emph{left} and \emph{right} are the distinguished children of
\defref{c01-binary}, heights stay in edges as in \chref{c01}, and a \textbf{visit} never
means ``print'' specifically -- it means whatever the problem does with a node.
